% =====================================================================
% Modelo ABNT de artigo cientifico
% EguaTeX (www.eguatex.com.br). Licenca CC0 1.0: dominio publico.
%
% Segue a ABNT NBR 6022:2018 e as normas que ela chama: NBR 6023:2025
% (referencias), NBR 6024:2012 (secoes), NBR 6028:2021 (resumo) e
% NBR 10520:2023 (citacoes).
%
% Se a revista tiver modelo proprio, siga o modelo da revista: a NBR 6022
% deixa o projeto grafico a criterio do editor.
%
% Como usar: troque cada texto de exemplo pelo seu conteudo, coloque as
% suas obras no arquivo referencias.bib e compile com pdfLaTeX e BibTeX
% (o EguaTeX faz isso sozinho). A formatacao fica no arquivo abnt-eguatex.sty.
% =====================================================================
\documentclass[12pt,a4paper,oneside]{article}
\usepackage[artigo]{abnt-eguatex}

% Nota de autor com chamada propria (asterisco), diferente das notas
% numericas do texto (NBR 6022:2018, 5.1.2).
\newcommand{\notadeautor}[2]{#1\begingroup\renewcommand{\thefootnote}{\fnsymbol{footnote}}%
  \footnote{#2}\addtocounter{footnote}{-1}\endgroup}

\begin{document}

% =====================================================================
% ELEMENTOS PRE-TEXTUAIS (NBR 6022:2018, 5.1)
% =====================================================================

% Titulo e subtitulo, separados por dois-pontos, no idioma do texto (5.1.1)
\begin{center}
  {\bfseries TÍTULO DO ARTIGO: subtítulo, se houver}\par
  \vspace{\baselineskip}
  % Titulo em outro idioma, opcional, logo abaixo do titulo
  {\itshape Title of the article: subtitle, if any}\par
\end{center}

% Autores em ordem direta: prenome e sobrenome (5.1.2). O curriculo
% resumido, a vinculacao e o contato ficam em nota de rodape.
\vspace{\baselineskip}
\begin{flushright}
  \setcounter{footnote}{0}
  \notadeautor{Nome do Primeiro Autor}{Titulação. Instituição de vínculo. E-mail: autor1@exemplo.org.br.}\par
  \setcounter{footnote}{1}
  \notadeautor{Nome do Segundo Autor}{Titulação. Instituição de vínculo. E-mail: autor2@exemplo.org.br.}\par
\end{flushright}
\setcounter{footnote}{0}

% Resumo no idioma do texto, de 100 a 250 palavras (NBR 6028:2021, 4.1.8)
\begin{resumo}
Aqui vai o resumo do artigo, em um único parágrafo, com frases curtas e o verbo na terceira pessoa. Ele informa o objetivo, o método, os resultados e as conclusões. Em artigos de periódicos, o resumo tem de 100 a 250 palavras.
\palavraschave{palavra-chave; segunda palavra-chave; terceira palavra-chave.}
\end{resumo}

% Resumo em outro idioma, opcional, depois do resumo no idioma do texto (5.1.3)
\begin{otherlanguage}{english}
\begin{resumo}[Abstract]
Here goes the abstract, the translation of the summary into another language. It follows the same rules as the summary in Portuguese.
\palavraschave[Keywords]{keyword; second keyword; third keyword.}
\end{resumo}
\end{otherlanguage}

% Datas de submissao e aprovacao: dia, mes e ano (obrigatorio, 5.1.4)
% Identificacao e disponibilidade: DOI ou endereco (opcional, 5.1.5)
\par\vspace{\baselineskip}\noindent
Data de submissão: dia mês ano.\\
Data de aprovação: dia mês ano.\\
DOI: \url{https://doi.org/10.0000/exemplo}

% =====================================================================
% ELEMENTOS TEXTUAIS (NBR 6022:2018, 5.2)
% O nome das secoes fica a seu criterio.
% =====================================================================

\section{Introdução}

A introdução apresenta o assunto do artigo, os objetivos da pesquisa e o que mais for preciso para situar o tema. O artigo usa letra tamanho 12 e espaçamento simples em todo o texto.

\section{Desenvolvimento}

O desenvolvimento expõe o assunto de forma ordenada e pode ser dividido em seções e subseções.

\subsection{Seção secundária}

As citações seguem o sistema autor-data, como em \cite[p.~15]{livro2024} ou \citeonline[p.~15]{livro2024}. No código, a citação é feita com \texttt{\textbackslash cite\{chave\}}, e a obra entra sozinha na lista de referências. Uma citação direta de até três linhas fica no parágrafo, entre aspas duplas: “aqui vai o trecho copiado da obra” \cite[p.~32]{livro2024}. Uma citação indireta não leva aspas \cite{artigo2025}.

\begin{citacao}
Aqui vai uma citação direta com mais de três linhas. Ela tem recuo de 4 cm da margem esquerda, letra menor que a do texto, espaçamento simples e não leva aspas. A fonte vem no final, entre parênteses, com o sobrenome, o ano e a página \cite[p.~47]{capitulo2023}.
\end{citacao}

Notas explicativas usam números\footnote{Aqui vai uma nota explicativa, com letra menor e espaçamento simples.}, e a chamada das notas de autor usa asterisco, para não se confundir com elas.

\subsection{Ilustrações e tabelas}

Toda ilustração é citada no texto, como a Figura~\ref{fig:exemplo}, com o título em cima e a fonte embaixo.

\begin{figure}[htbp]
  \begin{ilustracao}{0.7\textwidth}
    \caption{Título da figura}\label{fig:exemplo}
    \fbox{\parbox[c][4cm][c]{\dimexpr\linewidth-2\fboxsep-2\fboxrule}{\centering Coloque a imagem aqui com\\\texttt{\textbackslash includegraphics[width=\textbackslash linewidth]\{arquivo\}}}}
    \fonte{elaborada pelos autores (2026).}
  \end{ilustracao}
\end{figure}

As tabelas seguem as normas de apresentação tabular do IBGE, como a Tabela~\ref{tab:exemplo}.

\begin{table}[htbp]
  \begin{ilustracao}{0.7\textwidth}
    \caption{Título da tabela, com o que, onde e quando}\label{tab:exemplo}
    \centering
    \begin{tabular*}{\linewidth}{@{\extracolsep{\fill}}lrr}
      \toprule
      Categoria & Quantidade & Percentual (\%) \\
      \midrule
      Categoria A & 120 & 60,0 \\
      Categoria B & 80 & 40,0 \\
      \midrule
      Total & 200 & 100,0 \\
      \bottomrule
    \end{tabular*}
    \fonte{elaborada pelos autores (2026).}
  \end{ilustracao}
\end{table}

\section{Considerações finais}

As considerações finais respondem aos objetivos e às hipóteses apresentados na introdução.

% =====================================================================
% ELEMENTOS POS-TEXTUAIS (NBR 6022:2018, 5.3)
% Ordem: referencias, glossario, apendice, anexo e agradecimentos (o
% ultimo).
% =====================================================================
% Referencias (obrigatorias; NBR 6023:2025). A lista e montada sozinha
% a partir do arquivo referencias.bib: entram so as obras citadas no texto,
% em ordem alfabetica e no formato da norma.
\bibliography{referencias}

\apendice{Título do apêndice}
Aqui vai o conteúdo do apêndice, feito pelos autores. Apague se não houver.

\anexo{Título do anexo}
Aqui vai o conteúdo do anexo, de outra autoria. Apague se não houver.

\titulopostextual{Agradecimentos}
Aqui vão os agradecimentos, em texto curto, aprovado pela revista.

\end{document}
